\documentclass[conference]{support/IEEEtran}
\IEEEoverridecommandlockouts

% IEEE conference version of the ACT Workshop draft (same text and figure
% placeholders). Body sections live in sections/, references in references/,
% and the IEEEtran class/bibliography style in support/.
\usepackage{cite}
\usepackage{amsmath,amssymb,amsfonts}
\usepackage{graphicx}
\usepackage{textcomp}
\usepackage[table]{xcolor}
\usepackage{booktabs}
\usepackage{multirow}
\usepackage{tikz}
\usetikzlibrary{arrows.meta,positioning,calc,fit}
\usepackage[hidelinks]{hyperref}

\newcommand{\system}{\textsc{AS-Depth}}
\definecolor{ACMRed}{HTML}{BF0000}
\newcommand{\todo}[1]{\textcolor{ACMRed}{[TODO: #1]}}
\definecolor{paperblue}{HTML}{24577A}
\definecolor{paperlight}{HTML}{EEF3F7}
\definecolor{papergreen}{HTML}{347264}
% acmart's accessibility description has no IEEEtran counterpart.
\providecommand{\Description}[1]{}
\newcommand{\pending}{\textcolor{paperblue}{\textit{Pending}}}
\newcommand{\draftnote}[1]{\par\smallskip\noindent
  {\small\color{paperblue}\textit{Draft note: #1}}\par\smallskip}
\newcommand{\figslot}[3]{%
  \IfFileExists{#1}{\includegraphics[width=\linewidth,height=#2,keepaspectratio]{#1}}{%
  \fcolorbox{paperblue!45}{paperlight}{%
  \parbox[c][#2][c]{\dimexpr\linewidth-2\fboxsep-2\fboxrule\relax}{%
  \centering\color{paperblue}\small #3\par\medskip
  {\footnotesize Figure reserved for the advisor-review draft}}}}}
\newcommand{\resultslot}[2]{\fcolorbox{paperblue!45}{paperlight}{%
  \parbox[c][#1][c]{\dimexpr\linewidth-2\fboxsep-2\fboxrule\relax}{%
  \centering\small\color{paperblue}#2\par\medskip
  \footnotesize Measurements pending; no synthetic results shown.}}}

\begin{document}

\title{AS-Depth: Unleashing OFA Collaboration to Accelerate Metric Depth Completion on SoC}

\author{\IEEEauthorblockN{Anonymous Author(s)}
\IEEEauthorblockA{Anonymous Institution\\
Anonymous}}

\maketitle

\begin{abstract}
Depth completion reconstructs dense metric depth from sparse or noisy measurements captured by low-cost sensors,
offers great benefits to autonomous driving and robotics.
However, industrial robotic systems often benefit from real-time metric depth inference,
which requires the model to be deployed on a resource-constrained platform (e.g., SoC) with limited computing power and energy budget.
Moreover, existing depth completion models are computationally expensive,
thus most of the research on depth completion evaluates on data centers GPUs, while the targeted SoC has much smaller computation budget.
We observe that the backbone of most depht completion models have clear separation between the encoder and decoder stages, which naturally lends itself to heterogeneous acceleration by offloading one stage to a dedicated accelerator.Building on this insight, we propose as-depth, a heterogeneous system dedicated for accelerating depth completion on SoCs via three techniques: (1) a depth aware QAT strategy and accelerator implemented on FPGA to maintain model accuracy under low precision, (2)an asynchronous pipeline to overlap encoding and decoding stages (3)a heterogeneous execution framework that efficiently partitions the computation between the SoC and the FPGA.
Extensive experiments show that as-depth achieves end to end throughput of 30 fps while maintaining the near SOTA metric-depth accuracy.
Compared with baselines run on single-GPU, AS-Depth achieves a 5.4$\times$ speedup and xxx times improvement in energy efficiency.
Our method further demonstrates that a complete model need not reside on one processor:
a specific module can be offload to another device,
which provides a reusable path for accelerating compatible modules from other models.
\end{abstract}

\begin{IEEEkeywords}
architecture \& system design, depth completion, heterogeneous computing, SoC, real-time inference
\end{IEEEkeywords}

\begin{figure}[!h]
  \centering
  \figslot{figures/depth-completion-teaser.pdf}{1.18in}{%
    Application teaser and depth-completion pipeline\\[5pt]
    Raw RGB-D with missing/noisy depth $\rightarrow$ completed metric depth
    $\rightarrow$ robotic perception\\[3pt]
    Include a compact Thor--FPGA partition inset}
  \caption{Depth completion repairs missing and unreliable measurements from
  commodity RGB-D sensors before metric depth is consumed by downstream
  robotic perception. The inset previews the encoder--decoder partition used
  by \system.}
  \Description{A teaser showing a noisy RGB-D input, a completed metric-depth
  output, a downstream robotics use case, and a compact overview of the
  Thor--FPGA execution path.}
  \label{fig:application-teaser}
\end{figure}
\section{Introduction}
\label{sec:introduction}

Metric depth measures the distance of each pixel in an image from the camera, 
providing geometry cues such as spatial and shape information that the 2D color 
image cannot provide. 
Such depth information can be inputted into the backbone model of 
applications such as autonomous driving and robotics shown in Figure 1, 
which greatly improves
the capability of spatial reasoning and enables precision interaction 
with the objects.

\begin{table}[t]
  \caption{Depth-completion frame rates on the listed platforms. The
  30\,FPS entry for AS-Depth is a design target, not a measurement.}
  \label{tab:intro-framerate}
  \centering
  \small
  \setlength{\tabcolsep}{3pt}
  \renewcommand{\arraystretch}{1.2}
  \begin{tabular}{l|ccccc}
    \hline
    & CDM & PriorDA & PromptDA & OMNI-DC & Ours \\
    \hline
    Platform & RTX 4090 & RTX 4090 & RTX 4090 & RTX 3090 & SoC \\
    \hline
    FPS & 6.62 & 6.49 & 5.32 & 4.26 & 30 \\
    \hline
    Quality & - & - & - & - & - \\
    \hline
  \end{tabular}
\end{table}

With the widespread availability of commodity depth cameras, acquiring depth information appears to be low cost and straightforward. 
However, commodity RGB-D cameras frequently suffer from missing values, 
depth holes, flying pixels, and measurement noise, 
especially under challenging materials such as transparent, 
reflective, and specular surfaces. 
Although recent studies have integrated depth into robotic manipulations, performance remain limited by the poor quality of depth data produced by such device. Therefore, evaluations of robotics is typically restricted to simulators where clean depth is available. 
Thus, depth completion emerges as an effective method for generating dense metric depth from commodity depth cameras, and has become a critical component for bridging the Sim to Real gap. 

Recent studies 1234 have demonstrated advanced accuracy of depth completion, but their increasingly large Transformer-based architectures impose substantial computation costs. The selected GPU baselines in Table~\ref{tab:intro-framerate} range from 4.26 to 6.62 FPS on RTX 3090/4090 platforms. Because the models and platforms differ, these figures provide context rather than a controlled speedup comparison. Achieving real-time depth completion on an edge SoC such as NVIDIA Jetson AGX Thor remains challenging under its computation budget, motivating algorithm and hardware acceleration.

% To mitigate the contradiction of model heavy computation cost and SoC's computation budget, we propose a novel method: offloading one part of the existing model's computation to another device --- specifically, a FPGA --- to accelerate the whole depth completion pipeline. By leveraging the FPGA's flexible hardware acceleration capabilities, we aim to achieve real-time performance on edge platforms. We begin by profiling the depth completion pipeline and the operator type, and then analyzing the computational workload distribution between the SoC and the FPGA to find the optimal partitioning strategy. We also explore the best customized-precision strategies for the FPGA implementation, and maintain the model accuracy throughout quantization aware training. 

To address this, we propose As-Depth, a asynchronous heterogeneous acceleration system built on SoC and FPGA tailored for depth completion task on edge platforms.
As-Depth achieves efficient real time rendering speed and SOTA rendering quality on the edge platform. At the architecure level, it decouples the image representation encoding from the metric depth decoding stage, enabling these two tasks to be executed asynchronously on the SoC and FPGA according to their computational characteristics. On the algorithmn-hardware co-design level, our depth aware quantization and customized decoder design effectively preserve the accuracy of metric depth while reducing computational overhead. At the accelerator level, our depth aware accelerator integrates an ethernet interface unit and a depth decoding engine to handle the decoding pipeline, enabling efficient data communication and computation.

Building on these insights above, we develop a novel heterogeneous acceleration system, As-Depth, that combines the strengths of both the SoC and the FPGA to achieve real-time depth completion on edge platforms. It leverages the mature CUDA cores and Tensor Cores of the SoC to accelerate the transformer-based encoder, while offloading the CNN decoder to our customized FPGA implementation. This design ensures the maximum utilization of both the SoC and the FPGA, enableing pipelined execution to hide the latency of the offloaded operations. Through our quantization aware training and depth aware customized precision, our system strike a balance between performance and resource overhead. Our contributions are as follows:  

\begin{itemize}
  \item We propose a novel heterogenous system and combine task level asynchronous execution for depth completion based on the existing depth completion pipeline, efficiently mapping computing units according to the different computation characteristics of the SoC and FPGA, and using asynchronous optical queued communication to improve inference performance. 
  \item We design a depth-aware quantization strategy, and a customized precision scheme for the fpga implementation. Our approach successfully balance the accuracy of metric depth between model inference latency.
  \item We implement a prototype of our As-Depth system on the Nvidia Jetson AGX Thor and Intel Stratrix FPGA platform. Our prototype runs at real-time speed of 30FPS, demonstrating the feasibility of our heterogeneous acceleration approach.
  \item We implement and evaluate As-Depth on Synopsys. Experiments show that our system achieves 5.4x speedup and xxx improvement in energy efficiency over the latest SOTA implementation on the targeted SoC with the whole depth completion pipeline. To the best of our knowledge, this is the first work that achieves real-time 30FPS inference of depth completion on edge platforms.
  Our system enables real-time depth near sensor possible which contribute to mitigate the Sim to Real gap in the future rbotics applications. 
  Moreover, our approach demonstrates the potential for heterogeneous acceleration in resource-constrained environments, paving the way for future advancements in edge AI systems.
\end{itemize}


\begin{figure}[ht]
  \centering
  \figslot{figures/workload-characterization.pdf}{1.50in}{%
    Depth-completion workload characterization\\[5pt]
    (a) Thor-only stage latency and 33.3-ms budget\\
    (b) Encoder--decoder operator / compute-memory contrast\\
    Annotate the partition-boundary tensor size}
  \caption{Workload characterization of the selected depth-completion model.
  The final plot will report measured Thor-only stage latency and contrast the
  computational characteristics of the Transformer encoder and CNN decoder
  under one controlled inference configuration.}
  \Description{Two measured panels: a Thor-only latency breakdown relative to
  a 33.3 millisecond budget and a comparison of encoder and decoder workload
  characteristics, including the partition-boundary tensor size.}
  \label{fig:workload}
\end{figure}

\section{Background and Workload Characterization}
\label{sec:background}

\subsection{Depth Completion and the Encoder--Decoder Pipeline}

Depth completion estimates a dense metric-depth map from an RGB image and an
incomplete depth observation.  The input depth may contain holes, outliers, and
material-dependent noise, whereas the output must retain both absolute scale
and local geometry for downstream robotic perception.  Modern methods commonly
pair a representation encoder with a dense prediction decoder.  The encoder
extracts global visual and geometric features; the decoder fuses and upsamples
these features to the target depth resolution.

\system{} preserves this algorithmic decomposition and uses its boundary as a
system partition.  The selected depth-completion network and its trained
weights are the workload being accelerated, while this paper focuses on the
low-precision deployment, heterogeneous execution, and dedicated decoder
hardware.  The final manuscript will state the exact backbone, decoder,
input/output resolutions, and feature tensor exported at the partition.

\subsection{Computational Characteristics}

The two stages expose different opportunities for edge acceleration.  The
encoder contains Transformer blocks whose matrix operations and MLP layers can
use Thor's mature GPU and Tensor Core execution stack.  The decoder is composed
primarily of structured dense-prediction operators for which customized
parallelism, buffering, and precision can be implemented on FPGA.  This
difference motivates a stage-level mapping rather than repeatedly moving
individual operators between devices.

The argument must be grounded in measurements.  We therefore profile the
unmodified model on Thor and report the latency of preprocessing, encoder,
decoder, and postprocessing under a fixed resolution, precision, batch size,
and power mode.  We additionally characterize the dominant operator classes,
their arithmetic intensity or compute/memory demand, and the size of the
encoder--decoder boundary tensor.

\subsection{Edge SoC--FPGA Heterogeneous Execution}

An edge SoC offers a complete software environment and high-throughput tensor
units, but its compute and energy budget is shared with sensing, planning, and
control.  An FPGA provides configurable precision, dataflow, and memory
organization for a stable subgraph.  The combination is useful only when the
partition is coarse enough to limit communication and when execution across
devices is scheduled to hide, rather than merely add, transfer latency.

Prior GPU depth-completion systems establish the accuracy of large models, and
prior FPGA accelerators demonstrate the efficiency of specialized low-precision
datapaths.  \system{} connects these directions by retaining the encoder on the
SoC, specializing the decoder on FPGA, and treating the model boundary and
cross-frame schedule as first-class parts of the design.

\section{Motivation}
\label{sec:motivation}

Based on the workload characterization above, we identify two challenges in
deploying real-time depth completion on a heterogeneous edge platform.  Each
challenge exposes an opportunity that directly motivates one part of
\system's design.

\paragraph{Challenge 1: Monolithic execution misses the real-time budget.}

The Transformer encoder and CNN decoder have different operator mixes and
hardware preferences, yet a conventional deployment executes both stages
sequentially on one processor.  The resulting end-to-end latency includes the
full encoder, decoder, and runtime path for every frame.  Merely moving the
decoder to another device does not remove this serialization if communication
and accelerator execution remain on the critical path.

\paragraph{Opportunity 1: Task-level heterogeneous pipelining.}
This opportunity follows from the clear stage separation.
The clear encoder--decoder boundary permits a coarse-grained partition: the
Thor executes the encoder while the FPGA executes the decoder.  Across a
stream of frames, queued communication and double buffering can overlap the
encoder of frame $i+1$ with transfer and decoding of frame $i$.  This turns
device specialization into steady-state pipeline throughput rather than an
additional synchronous offload cost.

\paragraph{Challenge 2: End-to-end low precision can degrade metric depth.}

Low precision is beneficial on both sides of the partition.  In particular,
selected MLP computation in the Thor-side encoder and the FPGA-side decoder
can be executed at reduced precision.  A uniform quantization policy, however,
must simultaneously accommodate Transformer activation distributions,
depth-regression sensitivity, and the numerical contract at the cross-device
boundary.  Aggressive quantization can therefore save computation and
bandwidth while introducing scale errors or local depth artifacts.

\paragraph{Opportunity 2: Coordinated depth- and hardware-aware precision.}
This opportunity addresses the tension between efficiency and depth quality.
Rather than treating encoder quantization, feature transfer, and decoder
quantization as independent optimizations, \system{} coordinates their
precision choices.  Depth-aware quantization (DAQ) targets the layers and
signals that most affect metric-depth accuracy, while hardware-aware precision
choices map the encoder MLP and decoder datapath to the efficient modes of
Thor and FPGA, respectively.  The same boundary representation also reduces
the cost of cross-device feature transfer.

\begin{figure}[t]
  \centering
  \figslot{figures/motivation-evidence.pdf}{1.40in}{%
    Motivation evidence\\[5pt]
    (a) Sequential versus asynchronous frame schedule\\
    (b) Accuracy sensitivity of encoder-MLP and decoder quantization\\
    Use measured latency and depth error only}
  \caption{Evidence motivating \system. (a) A synchronous offload leaves
  communication and decoder execution on the critical path, whereas a
  task-level schedule overlaps work across frames. (b) Low-precision choices
  in both the encoder MLP and decoder must be coordinated to avoid excessive
  metric-depth error.}
  \Description{Two panels comparing sequential and asynchronous execution and
  showing measured depth-error sensitivity to low precision in the encoder
  MLP and decoder.}
  \label{fig:motivation}
\end{figure}

\section{AS-Depth Design}
\label{sec:design}

\system{} combines three mechanisms derived from the opportunities in
Section~\ref{sec:motivation}.  A task-level execution framework partitions the
model and overlaps consecutive frames; DAQ coordinates low precision across
the Thor encoder and FPGA decoder; and a dedicated decoder accelerator turns
the regular decoder workload into an efficient hardware dataflow.

\begin{figure}[t]
  \centering
  \figslot{figures/as-depth-system.pdf}{1.55in}{%
    AS-Depth system architecture\\[5pt]
    Thor: RGB-D preprocessing $\rightarrow$ encoder $\rightarrow$ quantized boundary feature\\
    Queues / double buffers / Ethernet $\leftrightarrow$ FPGA decoder\\
    Inset: cross-frame pipeline for frames $i$, $i+1$, and $i+2$}
  \caption{Overview of \system.  Thor executes the representation encoder and
  produces the cross-device feature tensor; the FPGA executes the depth
  decoder.  Queues and double buffers preserve frame order while overlapping
  encoding, communication, and decoding across consecutive frames.}
  \Description{A system block diagram of the Thor encoder, quantized feature
  boundary, communication queues, FPGA decoder, completed-depth output, and a
  cross-frame pipeline inset.}
  \label{fig:system-overview}
\end{figure}

\begin{figure}[t]
  \centering
  \figslot{figures/fpga-decoder-accelerator.pdf}{1.70in}{%
    Dedicated FPGA depth-decoder accelerator\\[5pt]
    Ethernet/interface engine $\rightarrow$ command scheduler $\rightarrow$ compute engines\\
    Weight/activation buffers $\leftrightarrow$ external memory\\
    Label weighted, residual, activation, and upsampling paths}
  \caption{Architecture of the dedicated FPGA depth-decoder accelerator.  The
  final figure will show the communication front end, programmable layer
  controller, weighted and auxiliary compute engines, local reuse buffers, and
  external-memory datapath.}
  \Description{A block diagram of the FPGA decoder accelerator and its
  communication, scheduling, compute, buffering, and memory components.}
  \label{fig:fpga-accelerator}
\end{figure}

\subsection{Task-Level Asynchronous Heterogeneous Execution Framework}
\label{sec:framework}

\paragraph{Model partition and tensor contract.}
The model is partitioned once at the encoder--decoder boundary.  Thor consumes
the RGB image and sensor depth, executes the encoder, and exports a quantized
feature tensor.  The FPGA consumes this tensor and returns the completed
metric-depth map.  The final paper will state the exact tensor shape, layout,
precision, scale metadata, and frame identifier used by this interface.

\paragraph{Queued cross-frame execution.}
The host runtime assigns separate buffers to encoder output, outbound
transfer, decoder input, and completed depth.  Producer--consumer queues make
ownership explicit and prevent a buffer from being reused before its consumer
finishes.  After pipeline fill, the encoder for the next frame can progress in
parallel with transfer and decoder execution for the current frame.  We report
both single-frame latency and steady-state throughput because only the latter
benefits fully from cross-frame overlap.

\paragraph{Communication and back pressure.}
The runtime batches feature metadata with the boundary tensor and preserves
frame order on the return path.  If either device becomes the bottleneck,
bounded queues apply back pressure instead of silently dropping or reordering
frames.  Section~\ref{sec:evaluation} measures transfer time separately from
useful computation so that communication hiding can be verified.

\subsection{Depth-Aware Quantization and Customized Precision}
\label{sec:daq}

DAQ treats quantization as an end-to-end property of the heterogeneous model.
On Thor, it targets selected MLP computation in the encoder; at the partition
boundary, it defines the representation transmitted to the FPGA; and inside
the FPGA, it determines the precision of decoder weights, activations,
accumulators, and output depth.  Training optimizes the quantized path against
metric-depth supervision so that numerical error is weighted by its effect on
the final depth map rather than by feature reconstruction alone, extending the
standard quantization-aware training setting~\cite{jacob2018quantization}.

\draftnote{Specify the encoder MLP layers covered by DAQ, the exact bit widths,
per-tensor or per-channel scaling, the depth-aware loss, calibration procedure,
and any layers retained at higher precision.  The ``without DAQ'' ablation must
keep the same partition and target precision while replacing DAQ with a
conventional quantization policy.}

\begin{figure}[t]
  \centering
  \figslot{figures/daq-precision.pdf}{1.42in}{%
    DAQ and precision map\\[5pt]
    Encoder MLP $\rightarrow$ boundary tensor $\rightarrow$ decoder\\
    Mark bit width, scale, accumulator, and depth-aware loss}
  \caption{Depth-aware quantization across the heterogeneous inference path.
  The final diagram will identify the quantized encoder MLP operations, the
  cross-device representation, decoder precision, and the depth-domain
  training signal.}
  \Description{A precision-flow diagram linking quantized encoder MLP layers,
  the transferred feature tensor, the FPGA decoder, and a depth-aware loss.}
  \label{fig:daq}
\end{figure}

\subsection{Dedicated FPGA Depth Decoder Accelerator}
\label{sec:accelerator}

The FPGA accelerator specializes the weighted and auxiliary operators in the
depth decoder while retaining a programmable layer schedule.  Its datapath
must expose the precision contract established by DAQ, reuse weights and
activations on chip where profitable, and stream larger intermediate tensors
through external memory without changing model semantics.  The communication
front end converts received feature packets into accelerator-visible buffers
and returns the completed depth map after the final layer.  This design follows
the broader principle of mapping quantized networks to customized FPGA
dataflows~\cite{umuroglu2017finn} while specializing the execution path for
metric-depth decoding.

\paragraph{Compute and dataflow.}
The architecture description will quantify parallel multiply--accumulate
lanes, accumulator precision, supported convolution variants, and the mapping
of activation, residual, and upsampling operators.  A layer controller selects
the appropriate engine and coordinates local buffers with external memory.

\paragraph{Control and correctness boundary.}
For every layer, the scheduler records tensor dimensions, addresses,
quantization parameters, and dependencies.  Verification compares each
hardware layer boundary and the final depth map against the quantized software
reference before end-to-end performance is reported.

\section{Evaluation}
\label{sec:evaluation}

We evaluate \system{} by answering four questions.  First, can heterogeneous
execution sustain the 30-FPS real-time target once communication and runtime
overheads are included?  Second, does DAQ preserve metric-depth quality on
challenging sensor data?  Third, how does the complete system compare with
recent depth-completion methods?  Finally, what FPGA resource, communication,
power, and energy costs are required to obtain these gains?

\subsection{Methodology}

\paragraph{Baselines.}
We use two groups of baselines for different purposes.  The controlled system
comparison runs the same depth-completion network in three deployments:
\emph{GPU-only} on an RTX~4090, \emph{Thor-only} on the target edge SoC, and
\system{} with its encoder on Thor and decoder on FPGA.  The GPU-only and
Thor-only baselines execute the complete model sequentially on one device;
\system{} includes feature transfer and output reconstruction in its
end-to-end time.  For the component study, \system{} without DAQ retains the
same partition, asynchronous runtime, accelerator, and target precisions, but
replaces the proposed depth-aware policy with conventional low-precision
quantization.  This makes the comparisons in Table~\ref{tab:methodology}
attributable to deployment and quantization choices rather than to different
networks.

For model-level context, we additionally compare against CDM, PriorDA,
PromptDA, and OMNI-DC, following the methods introduced in
Section~\ref{sec:introduction}.  We use their released checkpoints and official
evaluation procedures whenever available.  Because their published FPS values
come from different GPU platforms, the state-of-the-art throughput plot labels
the execution platform of every method and does not treat cross-platform FPS
as a controlled speedup claim.

\paragraph{Benchmarks.}
We evaluate all accuracy-oriented experiments on HAMMER and ClearPose.  These
datasets exercise different object geometries and difficult sensing conditions,
so results are reported for each dataset rather than averaged into one score.
Every method uses the same test split, validity mask, crop, input resolution,
and depth range within a dataset.  When a released method requires a different
input representation, only the documented preprocessing is changed; the
ground-truth mask and metric computation remain common.  The system and
ablation experiments use batch size one, which matches continuous RGB-D
capture on the target robotic platform.

\paragraph{Experiment platform.}
The prototype comprises an NVIDIA Jetson AGX Thor and an Intel Stratix FPGA
connected through the system interface described in Section~\ref{sec:framework}.
The FPGA clock, external-memory configuration, link rate, Thor power mode, and
software versions are held fixed across the two \system{} variants.  The
RTX~4090 baseline uses the same model weights, input resolution, preprocessing,
and output validation.  Before submission, the final platform table will state
the exact CUDA/TensorRT stack, FPGA device and achieved clock, host--FPGA link,
and precision assigned to every model stage.

\paragraph{Metrics and measurement protocol.}
We measure both single-frame latency and steady-state throughput.  Latency is
decomposed into preprocessing, encoder execution, boundary-tensor transfer,
FPGA decoding, return transfer, and output reconstruction; FPS is measured only
after the asynchronous pipeline is filled.  Timing includes synchronization
needed for a valid completed-depth output, rather than reporting kernel time
alone.  Each result will state the warm-up length, measured frame count, and
run-to-run variation.

Depth quality is evaluated with AbsRel, RMSE, MAE, and $\delta_{1.05}$.  Lower
is better for the three error metrics and higher is better for
$\delta_{1.05}$.  We report system power, energy per completed frame, and FPS/W
under a consistent measurement boundary.  FPGA resource results use post-fit
utilization and achieved frequency rather than pre-synthesis estimates.

\begin{table}[t]
  \caption{Configurations used for the controlled system comparison and
  component ablation.  ``Conventional LP'' and ``DAQ'' use the same target
  bit widths; they differ only in the quantization policy.}
  \label{tab:methodology}
  \centering
  \small
  \setlength{\tabcolsep}{3.0pt}
  \renewcommand{\arraystretch}{1.10}
  \begin{tabular}{lccc}
    \toprule
    Configuration & Encoder & Decoder & Precision policy \\
    \midrule
    RTX 4090 only & RTX 4090 & RTX 4090 & reference \\
    Thor only & Thor & Thor & reference \\
    \system{} w/o DAQ & Thor & FPGA & conventional LP \\
    \system{} w/ DAQ & Thor & FPGA & depth-aware LP \\
    \bottomrule
  \end{tabular}
\end{table}

\subsection{Ablation Experiments}

Figure~\ref{fig:ablation-results} evaluates the four configurations in
Table~\ref{tab:methodology}.  We interpret the bars cumulatively: the
single-device baselines establish the deployment cost, the system without DAQ
exposes the benefit of heterogeneous execution, and the complete system shows
whether depth-aware precision recovers quality without giving up the hardware
path.

\paragraph{Throughput and heterogeneous execution.}
Figure~\ref{fig:ablation-results}(a) reports absolute FPS and marks the 30-FPS
real-time threshold.  Thor-only versus \system{} without DAQ is the primary
controlled comparison for the Thor--FPGA partition: both process the same
frames with the same network, while only the latter offloads the decoder and
overlaps encoder, communication, and decoder work across frames.  We therefore
report not only their FPS difference, but also the stage-latency breakdown and
the fraction of transfer hidden by the asynchronous schedule.  \system{} with
and without DAQ is then compared under the same runtime to determine whether
the precision policy changes throughput.  GPU-only is retained as a high-end
performance reference, not as an edge-device energy baseline.  The completed
discussion will state absolute FPS first, followed by speedup over Thor-only,
and will explain any remaining gap using the slowest pipeline stage rather
than attributing it generically to ``offloading.''

\paragraph{Metric-depth accuracy.}
Figures~\ref{fig:ablation-results}(b)--(e) separate HAMMER and ClearPose and
show both an error metric and a threshold accuracy metric.  The decisive
comparison is \system{} without versus with DAQ because their model partition,
target bit widths, accelerator, and runtime are identical.  Any quality
difference therefore measures the value of coordinating the quantized encoder
MLP, boundary representation, and FPGA decoder using metric-depth supervision.
The GPU-only and Thor-only results provide the unpartitioned reference and
reveal the total degradation introduced by the heterogeneous low-precision
path.  We report each dataset independently and discuss whether improvements
come from reducing large depth errors, increasing the fraction of accurate
pixels, or both; a single aggregate ``near-SOTA'' statement is not used in
place of these measurements.

% The component ablation is the only full-width figure in the paper.
\begin{figure*}[t]
  \centering
  \figslot{figures/ablation-results.pdf}{1.62in}{%
    Results of ablation experiments --- grouped bar charts\\[5pt]
    AS-Depth w/ DAQ / AS-Depth w/o DAQ / GPU only / Thor only\\
    (a) FPS $\uparrow$ with 30-FPS line \quad
    (b) HAMMER AbsRel $\downarrow$ \quad
    (c) HAMMER $\delta_{1.05}$ $\uparrow$\\
    (d) ClearPose AbsRel $\downarrow$ \quad
    (e) ClearPose $\delta_{1.05}$ $\uparrow$}
  \caption{Results of the component ablation.  Throughput isolates the effect
  of heterogeneous execution, while the HAMMER and ClearPose panels quantify
  the accuracy recovered by DAQ under the same partition and target precision.}
  \Description{Five grouped bar charts comparing FPS, HAMMER absolute
  relative error and delta accuracy, and ClearPose absolute relative error and
  delta accuracy for AS-Depth with DAQ, AS-Depth without DAQ, GPU-only, and
  Thor-only.}
  \label{fig:ablation-results}
\end{figure*}

\subsection{Comparison with State-of-the-Art Methods}

\paragraph{Throughput.}
Figure~\ref{fig:sota-results}(a) compares the reported inference rate of CDM,
PriorDA, PromptDA, OMNI-DC, and \system{}.  Each bar is annotated with its
execution platform and the plot includes the 30-FPS reference line.  This
comparison answers a deployment question--whether current depth-completion
methods can service a real-time sensor stream--while the controlled
GPU-only/Thor-only comparison above supports architectural speedup claims.  In
the completed result discussion, we report \system{}'s measured end-to-end FPS
and identify whether its advantage comes from the coarse model partition,
cross-frame overlap, or low-precision execution, using the ablation evidence
rather than the cross-platform bars alone.

\paragraph{Depth accuracy.}
Figures~\ref{fig:sota-results}(b) and (c) compare all methods under common
HAMMER and ClearPose protocols.  The discussion reports the absolute metric
for every method, \system{}'s gap to the best-performing method, and its rank on
each dataset.  We then relate that gap to the deployment result: the relevant
question is whether \system{} reaches real-time edge throughput while keeping
the accuracy loss small and consistent across both datasets.  Claims of
state-of-the-art or near-state-of-the-art quality will be made only for metrics
and datasets for which the measured values support them.

\begin{figure}[t]
  \centering
  \figslot{figures/sota-results.pdf}{1.68in}{%
    Results of comparison with state of the art --- grouped bar charts\\[5pt]
    (a) FPS $\uparrow$: CDM / PriorDA / PromptDA / OMNI-DC / AS-Depth,\\
    with platform labels and 30-FPS line\\
    (b) HAMMER accuracy \qquad (c) ClearPose accuracy}
  \caption{Comparison with state-of-the-art depth-completion methods.
  Throughput is labeled by platform and provides deployment context; accuracy
  is evaluated on HAMMER and ClearPose with a common protocol.}
  \Description{Three grouped bar charts comparing FPS, HAMMER depth accuracy,
  and ClearPose depth accuracy across CDM, PriorDA, PromptDA, OMNI-DC, and
  AS-Depth.}
  \label{fig:sota-results}
\end{figure}

\subsection{Overhead Analysis}

\paragraph{FPGA resource and communication overhead.}
Table~\ref{tab:hardware-overhead} reports post-fit logic, register, DSP, and
on-chip-memory utilization together with achieved frequency.  To connect
resource cost to the system design, we also report boundary bytes per frame,
measured link latency, and standalone decoder latency.  These measurements
show whether the selected precision fits the target device and whether feature
transfer is sufficiently small to be hidden behind Thor encoding.  Resource
percentages are always reported with the device capacity and synthesis tool
settings so that raw counts are interpretable.

\paragraph{Power and energy efficiency.}
Power is measured for Thor-only, the active FPGA, and the complete
Thor--FPGA system using the same input stream and steady-state window used for
FPS.  We derive energy per frame as system power divided by completed frames
per second and additionally report FPS/W.  The analysis separates the extra
power introduced by the FPGA and link from the energy saved by shortening the
end-to-end execution time.  RTX~4090 power is reported as a separate reference
because its platform boundary differs from the embedded prototype.  The final
discussion will pair every energy-efficiency ratio with both its power and
throughput components, explaining when additional hardware power is offset by
a shorter end-to-end execution time.

\begin{table}[t]
  \caption{Post-fit hardware, communication, and energy overhead.  Values
  marked ``--'' await implementation or measurement.}
  \label{tab:hardware-overhead}
  \centering
  \small
  \setlength{\tabcolsep}{3.0pt}
  \renewcommand{\arraystretch}{1.08}
  \begin{tabular}{lr@{\qquad}lr}
    \toprule
    Metric & Value & Metric & Value \\
    \midrule
    Logic / registers & -- & DSP blocks & -- \\
    On-chip memory & -- & Frequency & -- \\
    Feature bytes/frame & -- & Link latency & -- \\
    Decoder latency & -- & System power & -- \\
    Energy/frame & -- & FPS/W & -- \\
    \bottomrule
  \end{tabular}
\end{table}

\section{Conclusion}
\label{sec:conclusion}

This paper presents \system, a heterogeneous depth-completion system that
decouples representation encoding on Thor from depth decoding on FPGA.  Its
task-level runtime overlaps consecutive frames, while DAQ coordinates
low-precision encoder MLP computation, the cross-device representation, and
the dedicated decoder datapath.  The final manuscript will close by reporting
the measured throughput, metric-depth accuracy, and energy efficiency achieved
by this combination on HAMMER and ClearPose.


\bibliographystyle{support/IEEEtran}
\bibliography{references/as-depth,references/architecture}
\end{document}
